\documentclass[10pt,letterpaper,sans]{moderncv}
\moderncvstyle{banking}
\moderncvcolor{blue}

\usepackage[scale=0.84]{geometry}
\nopagenumbers{}
\setlength{\hintscolumnwidth}{2.35cm}

\name{Krishnaa}{Vadivel}
\title{Software Engineer | Machine Learning | Distributed Computing}
\email{srikrishnaa.careers@gmail.com}
\phone[mobile]{516-853-5663}
\social[linkedin]{sri-krishnaa-ece-phd}
\extrainfo{\href{https://krishnaa423.github.io/industry-personal-website/}{Website} \textbar{} \href{https://github.com/krishnaa423}{GitHub: krishnaa423} \textbar{} \href{https://leetcode.com/u/krishnaa42342}{LeetCode: krishnaa42342} \textbar{} \href{https://hub.docker.com/u/krishnaa42342}{Docker Hub: krishnaa42342}}

\begin{document}
\makecvtitle

\section{Profile}
\cvitem{}{Software engineer and PhD researcher focused on machine learning, distributed computing, numerical software, and engineering tooling. Strong background in Python and C++ development for simulation, automation, GPU acceleration, data systems, and infrastructure spanning research and production-adjacent environments.}

\section{Research Experience}
\cventry{2021--present}{Research Assistant / PhD Researcher}{Yale University, Electrical Engineering}{}{}{\begin{itemize}%
\item Develop scientific software for first-principles and many-body workflows in materials research, combining Python and C++ with MPI, OpenMP, PETSc, SLEPc, CUDA, ROCm, and Linux-based automation.
\item Build distributed and accelerator-enabled workflows for HPC environments including Frontier, Frontera, and Perlmutter, with experience spanning NVIDIA and AMD GPU systems on national-lab clusters.
\item Create reusable tooling, notes, and structured workflows that bridge theory, implementation, debugging, and reproducibility in computational science.
\item Extend into ML-assisted and retrieval-driven automation where useful, particularly with LangChain, OpenAI API models, and local Ollama models for coding and research-workflow setup.
\item Co-authored a 2025 \textit{Journal of the American Chemical Society} paper and presented APS talks in 2024, 2025, and 2026 on self-trapped exciton formation and related first-principles work.
\end{itemize}}

\section{Professional Experience}
\cventry{2019--2021}{Software and Embedded Systems Engineer}{Probe Technologies Inc. / Weatherford Wireline Services}{}{}{\begin{itemize}%
\item Built CUDA-accelerated processing software for large acoustic waveform and sensor datasets used in engineering diagnostics and field-tool analysis.
\item Developed ASP.NET services and Microsoft SQL Server workflows for internal software, field-data management, and remote tool-calibration support.
\item Developed C\# GUI applications for preparing customer inputs and interacting with engineering workflows tied to deployed hardware systems.
\item Worked close to embedded and instrumentation systems, giving practical experience with software operating alongside deployed hardware and field constraints.
\end{itemize}}
\cventry{2018--2019}{Photonics Technical Writer}{FindLight}{}{}{\begin{itemize}%
\item Produced technical content for optics and photonics topics, demonstrating strong technical communication across specialized engineering subject matter.
\end{itemize}}

\section{Selected Projects}
\cventry{}{Equivariant Graph Neural Networks for Molecular Modeling}{}{}{}{\begin{itemize}%
\item Built PyTorch Geometric workflows for molecular-property prediction using equivariant graph neural networks and graph-transformer-style modeling ideas on QM9-style data.
\end{itemize}}
\cventry{}{Agentic Coding and Scientific Workflow Generation}{}{}{}{\begin{itemize}%
\item Developed LangChain-based retrieval-driven tooling using OpenAI API models and local Ollama models to read documentation and academic papers, then generate structured starter inputs and workflow artifacts.
\end{itemize}}
\cventry{}{Scientific Visualization and GPU Tooling}{}{}{}{\begin{itemize}%
\item Built CUDA-based tooling and related utilities for large numerical datasets, interactive inspection, and engineering communication.
\end{itemize}}
\cventry{}{Photonics and Simulation Code Artifacts}{}{}{}{\begin{itemize}%
\item Created clean, inspectable software artifacts for FDTD-based photonics simulation and related numerical-science workflows.
\end{itemize}}

\section{Education}
\cventry{2021--present}{PhD, Electrical Engineering}{Yale University}{}{}{Ongoing work combining scientific computing, high-performance workflows, and computational materials research.}
\cventry{2023}{MPhil, Electrical Engineering}{Yale University}{}{}{}
\cventry{2023}{MS, Electrical Engineering}{Yale University}{}{}{}
\cventry{2017--2018}{MEng, Electrical and Computer Engineering}{Cornell University}{}{}{}
\cventry{2013--2017}{BS, Electrical and Computer Engineering}{Cornell University}{}{}{}

\section{Selected Coursework}
\cvitem{Computing Foundations}{Theoretical Computer Science; Probabilistic Machine Learning}
\cvitem{Scientific Foundations}{Graduate Quantum Mechanics; Solid State Physics; Many-Body Quantum Physics}
\cvitem{Engineering Foundations}{Signals and systems, embedded systems, semiconductor physics, and numerical/scientific computing practice}

\section{Technical Skills}
\cvitem{Languages}{Python, C, C++, JavaScript, Bash, PowerShell, SQL}
\cvitem{Scientific Computing}{MPI, OpenMP, PETSc, SLEPc, mpi4py, petsc4py, slepc4py, NumPy, pandas, SciPy, SymPy, matplotlib}
\cvitem{Systems}{Linux command line, CMake, Docker, Docker Compose, Kubernetes, gcloud}
\cvitem{Platforms}{CUDA, ROCm, Frontier, Frontera, Perlmutter}
\cvitem{ML / Automation}{PyTorch, PyTorch Geometric, scikit-learn, LangChain, OpenAI API, Ollama, graph neural networks, graph transformers, agentic workflows}

\end{document}
